\documentclass[12pt,letterpaper]{article}

\usepackage[spanish]{babel}
\usepackage[utf8]{inputenc}
\usepackage[T1]{fontenc}
\usepackage{graphicx}
\usepackage{fancyhdr}
\usepackage{geometry}
\usepackage{amsmath, amssymb}
\usepackage{titlesec}

\geometry{margin=2.5cm}
\setlength{\headheight}{52pt}


\pagestyle{fancy}
\fancyhf{}

\fancyhead[L]{
    \includegraphics[height=1.5cm]{ur_image.png}
}

\fancyhead[R]{
\begin{tabular}{r}
\textbf{Monitoria Álgebra Lineal} \\
Gabriel Fernando Salcedo Zamora\\
\end{tabular}
}

\renewcommand{\headrulewidth}{0.4pt}

\begin{document}

\section*{Ejercicios factorización LU}
\begin{enumerate}
    \item Encuentre la matriz diagonal inferior $L$ y la matriz diagonal superior $U$ de las siguientes matrices:
    \begin{enumerate}
        \item \begin{pmatrix} 1 & 4 & 6\\2 & -1 & 3 \\ 3 & 2 & 5 \end{pmatrix}
        \item \begin{pmatrix} 2 & 3 & -1 & 6\\4 & 7 & 2 & 1 \\ -2 & 5 & -2 & 0 \\ 0 & -4 & 5 & 2 \end{pmatrix}
    \end{enumerate}
    \item Resuelva el sistema utilizando factorización $LU$. Recuerde que $Ax=LUx=b$
    \begin{enumerate}
        \item A = \begin{pmatrix} 1 & 4 & 6\\2 & -1 & 3 \\ 3 & 2 & 5 \end{pmatrix}
    \quad 
    b = \begin{pmatrix} -1 \\ 7 \\ 2\end{pmatrix},
        \item A = \begin{pmatrix} 3 & 9 & -2\\6 & -3 & 8 \\ 4 & 6 & 5 \end{pmatrix}
    \quad 
    b = \begin{pmatrix} 3 \\ 10 \\ 4\end{pmatrix},
    \end{enumerate}
\end{enumerate}

\newpage

\section*{Respuestas}
\begin{enumerate}
    \item 
    \begin{enumerate}
        \item \[
        L = \begin{pmatrix} 1 & 0 & 0\\2 & 1 & 0 \\ 3 & \frac{10}{9} & 1 \end{pmatrix}
        \quad 
         U = \begin{pmatrix} 1 & 4 & 6\\0 & -9 & -9 \\ 0 & 0 & -3 \end{pmatrix}\]
        \item \[
        L = \begin{pmatrix} 1 & 0 & 0 & 0 \\2 & 1 & 0 & 0 \\ -1 & 8 & 1 & 0 \\ 0 & -4 & -\frac{3}{5} & 1 \end{pmatrix}
        \quad 
        U = \begin{pmatrix} 2 & 3 & -1 & 6\\0 & 1 & 4 & -11 \\ 0 & 0 & -35 & 94 \\ 0 & 0 & 0 & \frac{72}{5} \end{pmatrix}\]
    \end{enumerate}
    \item 
    \begin{enumerate}
        \item x = \begin{pmatrix}
            -\frac{1}{3} \\[6pt] -\frac{8}{3} \\[6pt] \frac{5}{3} \end{pmatrix}
        \item x = \begin{pmatrix}
\frac{275}{63} \\[6pt]
-\frac{12}{7} \\[6pt]
-\frac{8}{3}
\end{pmatrix}
    \end{enumerate}
\end{enumerate}
\end{document}
